\begin{apartat}{1,25}
Justifiqueu si es podria determinar la massa d'un objecte penjant-lo d'una molla de constant elàstica coneguda (100~N/m) i deixant-lo oscil·lar unes quantes vegades i, en cas afirmatiu, expliqueu com la calcularíeu. Obtindríem el mateix resultat si ho féssim a la Lluna? Negligiu l'efecte de la força de fricció.
\end{apartat}

\begin{apartat}{1,25}
Deduïu l'equació de moviment de l'objecte a partir de l'equació del moviment harmònic simple (MHS) tenint en compte que l'amplitud del moviment és de 6,00~cm, que la freqüència d'oscil·lació és de 10,0~Hz i que el moviment s'inicia quan l'acceleració és màxima i positiva. Calculeu la velocitat i l'acceleració màximes del MHS a partir de l'equació de moviment.
\end{apartat}
